% Generated draft: requires mathematical and source review.
\documentclass[11pt]{article}
\usepackage[margin=25mm]{geometry}
\usepackage{fontspec}
\setmainfont{Latin Modern Roman}
\usepackage{amsmath,amssymb,mathtools,unicode-math}
\setmathfont{Latin Modern Math}
\usepackage{enumitem,needspace,xurl,hyperref}
\setlength{\parindent}{0pt}
\setlength{\parskip}{5pt}
\setlength{\emergencystretch}{3em}
\newcommand{\N}{\mathbb{N}}
\newcommand{\Z}{\mathbb{Z}}
\newcommand{\Q}{\mathbb{Q}}
\newcommand{\R}{\mathbb{R}}
\newcommand{\C}{\mathbb{C}}
\newcommand{\F}{\mathbb{F}}
\newcommand{\A}{\mathbb{A}}
\newcommand{\PP}{\mathbb{P}}
\newcommand{\E}{\mathbb{E}}
\newcommand{\eps}{\varepsilon}
\newcommand{\dd}{\,\mathrm{d}}
\begin{document}
\section*{UnsolvedMath ID: 801010}
\textbf{Source problem code:} catalogue entry
\subsection{Definitions and mathematical statement}
The supplied source asks three existence questions about maximal finitely periodic groups. It does not define ``finitely periodic,'' specify the ambient group, or identify the topology used for the descriptive complexity requirements. Consequently, the following is a conditional formalization of the questions, with those missing data left explicit.

Let \(S\) be an ambient group equipped with a topology \(\tau\), and let \(\mathcal F\) be the intended class of finitely periodic subgroups of \(S\). These are unresolved parameters, not definitions recovered from the source. For \(G\leq S\), define maximality within this class by
\[
\operatorname{Max}_{\mathcal F}(G)
\quad\Longleftrightarrow\quad
G\in\mathcal F\ \text{ and }\ 
\forall H\in\mathcal F\,
\bigl(G\subseteq H\Longrightarrow H=G\bigr).
\]
A subset \(G\subseteq S\) is \(K_\sigma\) if there are compact subsets \(K_n\subseteq(S,\tau)\), for \(n\in\N\), such that
\[
G=\bigcup_{n\in\N}K_n.
\]
Write \(\mathcal B(S,\tau)\) for the Borel sigma algebra generated by \(\tau\).

With the missing parameters fixed as intended by the source, the three questions are
\begin{align*}
&\exists G\leq S\quad
\bigl(\operatorname{Max}_{\mathcal F}(G)\ \land\ |G|\leq\aleph_0\bigr),\\
&\exists G\leq S\quad
\bigl(\operatorname{Max}_{\mathcal F}(G)\ \land\ G\text{ is }K_\sigma\bigr),\\
&\exists G\leq S\quad
\bigl(\operatorname{Max}_{\mathcal F}(G)\ \land\ G\in\mathcal B(S,\tau)\bigr).
\end{align*}
These are separate existence questions; the witnesses need not coincide. Here the first formula interprets ``countable'' as at most countable. Whether the intended setting excludes finite groups is not specified.

As motivation, the source attributes work on maximal cofinitary groups to Bart Kastermans, reference \([\mathrm{Kas}]\). It reports that maximal cofinitary groups cannot be \(K_\sigma\) or abelian, and that assuming the axiom of constructibility \(V=L\), a coanalytic maximal cofinitary group exists. These are supplied background assertions, not independently verified results here, and are not asserted for finitely periodic groups. The registry labels the problem open; this label is supplied curation rather than an independent mathematical assessment.
\subsection{Short English statement}
Does there exist a maximal finitely periodic group that is countable? Does one exist that is \(K_\sigma\), or one that is Borel? The source leaves the meaning of finitely periodic, the ambient group, and the relevant topology unspecified, so the displayed formulas retain these as unresolved parameters.
\subsection{Sources}
\begin{itemize}
\item Local UnsolvedMath registry, permanent ID \(801010\), exact source problem code \texttt{catalogue entry}; title: \emph{Open Problems in Mathematical Logic --- Problem 22}. The record attributes the problem to Barmpalias, Gao et al., with proposed year 2026, and supplies the status ``open.''
\item Supplied source list: Barmpalias, Gao et al., \emph{Open Problems in Mathematical Logic} (2026), Problem 22. Supplied URL: \url{https://arxiv.org/abs/2608.26628}. The record reports access on 2026-10-16 and extraction by \texttt{arxiv\_latex}; these are supplied provenance claims, not independently verified access or source checks.
\item The supplied statement cites Bart Kastermans as \([\mathrm{Kas}]\) for work on maximal cofinitary groups. No full bibliographic entry or URL for this reference is supplied.
\end{itemize}
\end{document}
